\section*{Capítulo \ref{ch:g6:earth-a-planet} --- La Tierra, un planeta del sistema solar}

\begin{solution}{exo:g6:earth-a-planet:1}
\underline{Mercurio}, \underline{Venus}, \underline{la Tierra},
\underline{Marte}, Júpiter, Saturno, Urano, Neptuno. Los \omterm{def:g6:earth-a-planet:planet-types}{planetas rocosos}
son los cuatro primeros.
\end{solution}

\begin{solution}{exo:g6:earth-a-planet:2}
``Un \omterm{def:g6:earth-a-planet:planet-types}{planeta rocoso} es pequeño y denso, con una superficie sólida. Un
\omterm{def:g6:earth-a-planet:planet-types}{planeta gigante} es muy grande, de densidad baja, hecho sobre todo de gases
y de líquidos''.
\end{solution}

\begin{solution}{exo:g6:earth-a-planet:3}
Venus (12\,104~km, frente a 12\,756~km de la Tierra). El más grande es
Júpiter (142\,984~km); el más pequeño es Mercurio (4\,879~km).
\end{solution}

\begin{solution}{exo:g6:earth-a-planet:4}
(a) Falso: Venus no tiene agua líquida; hace demasiado calor. (b) Verdadero
($-65$~$^\circ$C de media, frente a 15~$^\circ$C). (c) Falso: los cuatro
planetas más cercanos al Sol son los \omterm{def:g6:earth-a-planet:planet-types}{planetas rocosos}; los \omterm{def:g6:earth-a-planet:planet-types}{planetas
gigantes} son los cuatro más lejanos. (d) Verdadero: unos 0.7~gramos por
centímetro cúbico.
\end{solution}

\begin{solution}{exo:g6:earth-a-planet:5}
$4879 \div 12756 \approx 0.382$, es decir, el 38~\%: Mercurio es un poco
más de un tercio de ancho que la Tierra.
\end{solution}

\begin{solution}{exo:g6:earth-a-planet:6}
$120536 \div 12756 \approx 9.4$: unas 9 Tierras. Saturno es un \omterm{def:g6:earth-a-planet:planet-types}{planeta gigante}.
\end{solution}

\begin{solution}{exo:g6:earth-a-planet:7}
El planeta~A es un \omterm{def:g6:earth-a-planet:planet-types}{planeta gigante}: es muy grande y su densidad, de
1.1~gramos por centímetro cúbico, es baja, como la de los \omterm{def:g6:earth-a-planet:planet-types}{planetas gigantes}
del sistema solar (entre 0.7 y 1.7). El planeta~B es un \omterm{def:g6:earth-a-planet:planet-types}{planeta rocoso}: es
pequeño y su densidad, de 5.0, es alta, como la de los \omterm{def:g6:earth-a-planet:planet-types}{planetas rocosos}
(entre 3.9 y 5.5).
\end{solution}

\begin{solution}{exo:g6:earth-a-planet:8}
(a) $464 - 15 = 449$~grados más caliente. (b) $15 - (-65) = 80$~grados más frío.
\end{solution}

\begin{solution}{exo:g6:earth-a-planet:9}
Venus tiene una atmósfera muy espesa de dióxido de carbono, que atrapa el
calor del Sol como una manta gruesa y no lo deja escapar. Mercurio, más
cerca del Sol, casi no tiene aire, así que no puede conservar su calor.
\end{solution}

\begin{solution}{exo:g6:earth-a-planet:10}
$228 - 150 = 78$~millones de kilómetros más lejos. $78 \div 150 = 0.52$: la
distancia de Marte es un 52~\% mayor que la de la Tierra.
\end{solution}

\begin{solution}{exo:g6:earth-a-planet:11}
Hace miles de millones de años, Marte tenía \omterm{def:g1:landscapes-around-us:water}{ríos}, \omterm{def:g1:landscapes-around-us:water}{lagos} e \omterm{def:g5:water-shapes-the-land:flood}{inundaciones}: era
más cálido y más húmedo. Hoy su agua está helada, en el hielo de los
casquetes polares y bajo el suelo cerca de los \omterm{def:g3:continents-and-oceans:poles}{polos}. El agua líquida no
puede durar ahora en su superficie porque Marte es muy frío
($-65$~$^\circ$C de media) y porque su aire es tan tenue que el agua se
congela enseguida o se convierte en vapor.
\end{solution}

\begin{solution}{exo:g6:earth-a-planet:12}
(a) $8.9 \div 9.8 \approx 0.908$, es decir, el 91~\%. (b) Como un gemelo:
casi el mismo tamaño (el 95~\% del diámetro de la Tierra), casi la misma
atracción de la gravedad, y también es un \omterm{def:g6:earth-a-planet:planet-types}{planeta rocoso}. No como un
gemelo: es mucho más caliente (464~$^\circ$C frente a 15~$^\circ$C), su
aire es muchísimo más espeso y está hecho de dióxido de carbono con \omterm{def:g4:clouds-rain-snow:cloud}{nubes}
de ácido, y no tiene agua líquida.
\end{solution}

\begin{solution}{pb:g6:earth-a-planet:1}
\textbf{1.} Venus (108~millones de~km), la Tierra (150), Marte (228).

\textbf{2.} $150 - 108 = 42$: Venus está 42~millones de kilómetros más
cerca. $228 - 150 = 78$: Marte está 78~millones de kilómetros más lejos.

\textbf{3.} $6792 \div 12756 \approx 0.532$, es decir, el 53~\%: Marte es
más o menos la mitad de ancho que la Tierra.

\textbf{4.} Los tres son \omterm{def:g6:earth-a-planet:planet-types}{planetas rocosos}: sus densidades (5.2, 5.5 y 3.9
gramos por centímetro cúbico) son altas, y sus diámetros son pequeños
comparados con los de los \omterm{def:g6:earth-a-planet:planet-types}{planetas gigantes}.

\textbf{5.} El aire más espeso está en Venus (92~bares); el más tenue, en
Marte (0.01~bar).

\textbf{6.} $1 \div 0.01 = 100$: el aire de Marte aprieta cien veces menos
que el aire de la Tierra.

\textbf{7.} Su aire espeso de dióxido de carbono atrapa el calor del Sol y
no lo deja escapar, como una manta gruesa.

\textbf{8.} $92 \div 1 = 92$: en el suelo de Venus, el aire aprieta
\textbf{92 veces} más fuerte que en el suelo de la Tierra. Una sonda tiene
que tener paredes muy resistentes para no quedar aplastada (y también tiene
que resistir el calor y las \omterm{def:g4:clouds-rain-snow:cloud}{nubes} de ácido).

\textbf{9.} La temperatura media de la Tierra, 15~$^\circ$C, está entre
0~$^\circ$C y 100~$^\circ$C: el agua no está ni helada ni hirviendo, así
que sigue líquida.

\textbf{10.} En Venus, a 464~$^\circ$C, muy por encima de 100~$^\circ$C,
el agua herviría enseguida y se convertiría en vapor. En Marte, a
$-65$~$^\circ$C, muy por debajo de 0~$^\circ$C, se congelaría (y poco a
poco se convertiría en vapor, porque el aire es muy tenue).

\textbf{11.} Hace mucho tiempo, Marte era más cálido y más húmedo que hoy:
el agua líquida podía correr por su superficie. Su \omterm{def:g5:weather-and-climate:climate}{clima} ha cambiado
mucho: ahora es un \omterm{def:g5:weather-and-climate:desert}{desierto} frío y seco.

\textbf{12.} Por ejemplo: ``Queridas sondas: aquí, en la Tierra, tenemos
agua líquida, en \omterm{def:g2:water-on-earth:ocean}{océanos}, \omterm{def:g1:landscapes-around-us:water}{ríos} y \omterm{def:g1:landscapes-around-us:water}{lagos}. Nuestras temperaturas son suaves,
15~$^\circ$C de media, entre la congelación y la ebullición del agua.
Estamos a la distancia adecuada del Sol, ni demasiado cerca ni demasiado
lejos. Y nuestra atmósfera nos mantiene calientes, da a los seres vivos los
gases que necesitan y permite que el agua siga líquida. Por eso aquí es
posible la vida''.
\end{solution}
